% Full title: Corrected Representation Identity
\section{Corrected representation identity}
\label{app:corrected-identity}

The wavenumber assignment in the complementary-field argument is recorded
here to prevent a common mismatch.  If $w^{\mathrm e}$ solves the background
equation $(\Delta+k^2)w^{\mathrm e}=0$ and $v^{\mathrm i}$ is a solution of
the complementary interior-$k$ equation, the quasiperiodic extinction pair is
\begin{equation}
  -S^{\mathrm{qp}}_{k,\beta}
  [\gamma_N^{\mathrm i}v^{\mathrm i}]
  +D^{\mathrm{qp}}_{k,\beta}
  [\gamma_D^{\mathrm i}v^{\mathrm i}].
  \label{eq:corrected-exterior-pair}
\end{equation}
Both operators in \eqref{eq:corrected-exterior-pair} use the background
wavenumber $k$.  The inclusion representation uses the free-space pair
\begin{equation}
  -S_{k^0}[\gamma_N^{\mathrm i}w^{\mathrm i}]
  +D_{k^0}[\gamma_D^{\mathrm i}w^{\mathrm i}],
  \qquad k^0=n^0k.
  \label{eq:corrected-interior-pair}
\end{equation}
Mixing $D^{\mathrm{qp}}_{k^0,\beta}$ into
\eqref{eq:corrected-exterior-pair} invalidates the exterior $k$-Helmholtz
extinction identity.  Equations \eqref{eq:corrected-exterior-pair} and
\eqref{eq:corrected-interior-pair} agree with the wavenumber separation used
in \eqref{eq:center-reconstruction} and with the architecture of
\cite[Appendix~A]{BarnettGreengard2010}.
